% =====================================================================
% GeoVac: A Field Guide
%
% An identity statement for the Geometric Vacuum (GeoVac) framework.
% Audience: readers who want to understand the project's arc from
% origin to current frontier without first learning the machinery
% (Connes distance, propinquity convergence, Glanois cyclotomic basis,
% etc.).  This is the "what is GeoVac" document — narrative, sober,
% honest about scope.  Math is layered in where it's load-bearing for
% the story; deep technical content is cross-referenced to the
% appropriate paper rather than recapitulated.
%
% Status: drafted June 2026 after the Periods-reframe arc (v3.46.0)
% and the seven A-tier register closures of 2026-06-03 (Sprints A1
% scoping + A2 + A4 + A5 + A6 + A7 + A8 + A1-Matsubara).
% =====================================================================
\documentclass[11pt,a4paper]{article}

\usepackage[T1]{fontenc}
\usepackage[utf8]{inputenc}
\usepackage{lmodern}
\usepackage{amsmath, amssymb, amsthm, mathrsfs, mathtools}
\usepackage{geometry}
\geometry{margin=1in}
\usepackage{hyperref}
\hypersetup{colorlinks=true, linkcolor=blue!50!black, citecolor=blue!50!black, urlcolor=blue!50!black}
\usepackage[expansion=false]{microtype}
\usepackage{booktabs}
\usepackage{xcolor}

\newcommand{\sthree}{S^{3}}
\newcommand{\sfive}{S^{5}}
\newcommand{\R}{\mathbb{R}}
\newcommand{\C}{\mathbb{C}}
\newcommand{\N}{\mathbb{N}}
\newcommand{\Z}{\mathbb{Z}}
\newcommand{\Q}{\mathbb{Q}}
\newcommand{\Tr}{\mathrm{Tr}}
\newcommand{\Vol}{\mathrm{Vol}}
\newcommand{\MT}{\mathcal{MT}}

\title{GeoVac:\ A Field Guide \\
\large From a Packing Puzzle to Periods of the Spectral Triple}
\author{J.~Loutey\thanks{Independent researcher.
\texttt{jloutey@gmail.com}.  Developed in an AI-augmented agentic
workflow with Anthropic's Claude;\ all theorems, design choices, and
editorial decisions are the author's responsibility.}}
\date{June 2026}

\begin{document}
\maketitle

\begin{abstract}
The Geometric Vacuum (GeoVac) framework began as a packing puzzle that
produced something resembling the periodic table.  Over several years
of investigation, it acquired a precise structural identity:\ a
discrete spectral triple that is an instance of the Marcolli--van
Suijlekom 2014 graph-network construction at the data level and
extends it via a forced-graph claim and a master Mellin engine,
whose transcendentals classify as cyclotomic mixed-Tate periods at
level $\le 4$, and which separates --- through theorems,
not preference --- what the underlying geometry forces from what
empirical measurement supplies.  This field guide is the project's
identity statement.  It tells the arc from origin to current frontier
without recapitulating the technical machinery, and points each kind
of reader toward the right entry into the corpus.  The one-sentence
reading:\ \emph{GeoVac is a discrete spectral triple whose periods
classify the projective content of physics, sharply separating what
geometry forces from what measurement supplies.}
\end{abstract}

\tableofcontents

% =====================================================================
\section{How This Started}
\label{sec:origin}
% =====================================================================

A short note on the path that led to the framework, told in the first
person.  The reader can skip it without missing technical content;\ the
arc proper begins in \S\ref{sec:packing}.

In high school I looked at the periodic table and wondered why it had
that shape --- in particular, why the rows didn't grow uniformly but
in a stepped pattern.  I sketched what I thought might be a more
natural picture:\ a circular layout with the first row at the centre
and successive rows fanning outward.  I had no mathematics to back it
up;\ it was a visual hunch.

In college I learned summation notation and the discipline of
inductive proof, and I felt --- for the first time --- that I had
permission to do real mathematics.  I went back to the periodic
table.  The row lengths grew by a recurring $+4$ between successive
pairs.  I built a nested summation that captured the pattern and
arrived at $2n^2$.  I looked up the number.  It was the degeneracy of
the Schr\"odinger principal quantum level.  I felt something between
satisfaction and disappointment:\ the pattern was already known ---
to Schr\"odinger, who had derived it by solving a differential
equation on $\R^3$.  My combinatorial route to the same number was a
curiosity, not a discovery.  I set the question down.

Years later, listening to a guest on Sean Carroll's podcast, I heard
the suggestion that perhaps, instead of quantising gravity, one
should try to quantise spacetime --- a discrete substrate as
foundation, not as a tool for approximating something continuous.
The framing revived the old question.  This time I posed it in polar
coordinates and treated Planck's constant as a unit volume to be
packed.  I wrote down a small set of packing axioms.  The
construction, executed honestly, produced circles whose node counts
reproduced the periodic table.

\begin{quote}
\itshape
This struck me as deeply important.
\end{quote}

I wrote a first paper that I would now describe as a book report:\ a
narrative explanation of how the packing maps to the periodic table,
without any of the spectral or geometric machinery that came later.
That paper became Paper~0~\cite{loutey_paper0} in its mature form.  I
then set out to see how the nodes might relate to quantum-number
counting.  Everything in the rest of this field guide --- the
Fock~1935 identification, the natural-geometry hierarchy, the
spectral-triple rotation, the periods classification, the
forced-versus-free seam --- is what happened next.

What follows is the technical arc.

% =====================================================================
\section{The Packing Puzzle}
\label{sec:packing}
% =====================================================================

The framework's origin is a geometric construction:\ pack
Planck's constant.  Treat $\hbar$ as a unit volume to be filled, and
ask which node arrangements complete the packing without gaps.  In two
dimensions the node counts give $(l, m)$:\ the integer pair indexing
azimuthal and magnetic quantum numbers, before any quantum-mechanical
input.  In three dimensions, an inductive argument over periodic-table
row lengths fixes the principal and spin quantum numbers
$(n, s)$~\cite{loutey_paper0}.  The resulting graph has nodes labeled
by quantum-number tuples $(n, l, m, s)$ and edges encoding spectral
selection rules.

The graph Laplacian on this object has eigenvalues
\begin{equation}
\label{eq:graph_eigs}
\lambda_n = -(n^2 - 1), \qquad n \in \N_{\ge 1},
\end{equation}
with multiplicity $n^2$ on the unit-radius rescaling --- the spectrum
of the hydrogen atom, to within an overall scale.  The shape of the
periodic table fell out of the packing.

This was the seed:\ a combinatorial object that, when interrogated,
returned the same numerology as observed atomic structure.  Whether
that was a coincidence or a clue was the question the rest of the
project tried to answer.

% =====================================================================
\section{Fock 1935:\ Discrete and Continuous Are the Same Thing}
\label{sec:fock}
% =====================================================================

In 1935 Vladimir Fock observed~\cite{fock1935} that the hydrogen atom
on $\R^3$ has a hidden $SO(4)$ symmetry, and that the stereographic
projection $\R^3 \hookrightarrow \sthree$ converts the Schr\"odinger
equation into an angular eigenvalue problem on the unit three-sphere.
What hydrogen really is, from this point of view, is the
Laplace--Beltrami operator on the unit $\sthree$;\ the $1/r$ Coulomb
potential is the coordinate distortion produced by the projection
(specifically, the chordal-distance identity that takes the round
metric to the stereographic one).

The decisive observation for GeoVac is this:\ the graph Laplacian on
the packed lattice of \S\ref{sec:packing} is conformally equivalent to
the continuous $\sthree$ Laplace--Beltrami operator, in a precise sense
that can be checked symbolically.  Paper~7 of the GeoVac
corpus~\cite{loutey_paper7} provides $18$ independent symbolic proofs
of this equivalence (\texttt{tests/test\_fock\_projection.py} and
\texttt{tests/test\_fock\_laplacian.py}).  The universal scaling
constant relating the two is $\kappa = -1/16$, derivable from the
packing construction~\cite{loutey_paper0}.

This is the framework's load-bearing identification:
\begin{quote}
\itshape
The discrete Fock graph and the continuous unit-$\sthree$
Laplace--Beltrami operator are not two models of the hydrogen atom.
They are the same object in two coordinates.
\end{quote}
The implication is operational:\ where the continuous formulation
requires integration over $\R^3$ with a $1/r$ singularity, the
discrete formulation requires a sparse matrix eigenvalue problem with
integer entries.  Same physics, different cost.

% =====================================================================
\section{Natural Geometry:\ a Hierarchy}
\label{sec:natural_geometry}
% =====================================================================

Section~\ref{sec:fock}'s identification works because $\sthree$ is the
natural geometry for the hydrogen atom:\ it is the coordinate system in
which the Schr\"odinger equation separates fully into angular
variables.  For other systems --- molecules, multi-electron atoms,
polyatomics --- the natural geometry shifts.  The framework's chemistry
investigation (2024--2026, Phase~2 of the project) systematically
catalogued the natural geometries:

\begin{center}
\begin{tabular}{@{}cllll@{}}
\toprule
Level & System & Natural geometry & Best accuracy & Paper \\
\midrule
1 & H (1-center, 1e) & $\sthree$ (Fock) & $< 0.1\%$ & 7 \\
2 & H$_2^+$ (2-center, 1e) & Prolate spheroid & $0.0002\%$ & 11 \\
3 & He (1-center, 2e) & Hyperspherical & $0.019\%$ & 13 \\
4 & H$_2$ (2-center, 2e) & Mol-frame hyperspherical & $96.0\%$ $D_e$ & 15 \\
5 & LiH (composed) & Composed (Level 3 $\otimes$ Level 4) & $5.3\%$ $R_{\mathrm{eq}}$ & 17 \\
\bottomrule
\end{tabular}
\end{center}

The choice of geometry is determined by where separation of variables
occurs;\ in each case the discrete graph Laplacian on the natural
geometry reproduces the continuous Schr\"odinger problem with the
same conformal-equivalence identification as in
\S\ref{sec:fock}.

Phase~2 reached known structural ceilings.  Heavy molecules
($Z \ge 20$ or strong multi-determinant correlation) escape the
framework's accuracy targets through identifiable channels (the
``multi-focal-composition wall'' documented in six independent
sprints).  The framework's most natural quantum-simulation
application --- $O(Q^{2.5})$ Pauli scaling on the composed
architecture, $51\times$--$1712\times$ Pauli-count advantage over
published Gaussian baselines, $28$ molecules in the
library~\cite{loutey_paper14} --- shipped at the end of Phase~2.

But a different question began to dominate the conversation around the
project, and that question is where the rest of this guide goes.

% =====================================================================
\section{The Rotation:\ from Solver to Spectral Triple}
\label{sec:rotation}
% =====================================================================

Phase~3 was planned as quantum simulation.  Instead, the project
rotated to a different question:\ \emph{what is this thing we have
been working with?}

The graph Laplacian on $\sthree$, the conformal equivalence with the
continuous problem, the natural-geometry hierarchy that produces
discrete substrates for every electron configuration --- these are
collectively the data of a Riemannian spectral
geometry.  Made precise, that means an algebra of functions, a
Hilbert space of states, and a Dirac-type operator, all assembled
into a structure that mathematicians call a spectral triple
\cite{connes1994}.

The precise structural claim:

\begin{quote}
\itshape
GeoVac is a discrete spectral triple that is an \emph{instance} of
the Marcolli--van Suijlekom 2014 gauge-network construction
\cite{marcolli_vs2014} at the data level (refined by the
Perez-Sanchez 2024/2025 clarification that the continuum limit is
Yang--Mills without Higgs in general) and an \emph{extension} of it
at the structural level via the forced-graph claim (Bertrand
$\times$ Fock projection), the master Mellin engine, and the
cosmic-Galois layer $U^*$ (Paper~56).
\end{quote}

This claim was checked, not assumed.  Working hypothesis WH1 (in the
project's internal register, CLAUDE.md~\S~1.7) was elevated to
PROVEN status in May 2026 via a five-lemma Gromov--Hausdorff
convergence theorem on truncated Camporesi--Higuchi spectral
triples~\cite{loutey_paper38} (Paper~38, released to Zenodo
2026-05-07).  The proof rests on the Latr\'emoli\`ere propinquity
framework~\cite{latremoliere2017} for quantifying when two spectral
triples are close in the Connes--Hausdorff sense, with the asymptotic
rate constant
\begin{equation}
\label{eq:propinquity_rate}
\Lambda_{n_{\max}} \le C_3 \cdot \gamma_{n_{\max}},
\qquad \gamma_{n_{\max}} \xrightarrow{n_{\max} \to \infty}
\frac{4}{\pi},
\end{equation}
universal across compact connected Lie groups~\cite{loutey_paper40}.
The rate constant $4/\pi$ is the substrate's $\mathrm{Vol}(S^2) /
\pi^2$ signature --- the Hopf-base measure entering through a single
geometric mechanism, classified in
\S\ref{sec:periods} below.

Following Paper~38, the project produced eleven math.OA standalone
papers (Papers~38--49) extending the framework along the Lorentzian
axis~\cite{loutey_paper42, loutey_paper43, loutey_paper44,
loutey_paper45, loutey_paper46, loutey_paper47, loutey_paper48,
loutey_paper49}.  Sub-results include the first published Lorentzian
propinquity convergence theorem on truncated Krein spectral triples
(Paper~45)~\cite{loutey_paper45}, an explicit strong-form extension
(Paper~46)~\cite{loutey_paper46}, and the operator-system Lorentzian
pre-length space (OSLPLS) category that closes a published-open
question of Connes--van Suijlekom on Lorentzian extensions
(Paper~49)~\cite{loutey_paper49}.  These results were checked
independently against the four published-open questions Connes--van
Suijlekom 2021~\cite{connes_vs2021} explicitly deferred to ``elsewhere.''

\paragraph{Why this rotation mattered.}
The shift from ``how accurately can we solve quantum chemistry?'' to
``what is GeoVac?'' produced the right kind of answers.  Chemistry
accuracy is a sliding scale --- always tighter at more cost.  But
the spectral-triple structural reading is binary:\ either GeoVac IS
the data of a spectral triple, or it isn't.  Theorem-grade closure
is available either way.  The framework's identity is now in the
first category.

% =====================================================================
\section{Periods:\ Where Transcendentals Come From}
\label{sec:periods}
% =====================================================================

A spectral triple has a heat kernel, and the heat kernel produces
transcendentals.  In GeoVac, those transcendentals are not random.
They sit in named classes, with named mechanisms producing them.

\paragraph{The graph itself is $\pi$-free.}
At what the project calls \emph{Layer 1} --- the bare combinatorial
graph and its integer eigenvalue spectrum --- there are no
transcendentals.  Matrix elements are rationals (or, at most,
algebraic over $\Q$).  This has been checked across the corpus:
verified explicitly on the unit $\sthree$ Coulomb spectrum (Paper~7),
on the Hopf bundle gauge structure (Paper~25), on the Bargmann--Segal
$\sfive$ Hardy lattice (Paper~24), and on $28$ molecular
species in the qubit Hamiltonian library (Paper~14).

\paragraph{Transcendentals enter at projection.}
Layer~1 by itself does not produce physical observables.  Observables
come from \emph{Layer~2:}\ projections of the discrete substrate onto
continuous manifolds, integration regions, or measurement-coordinate
choices.  Each projection introduces a controlled transcendental
through a controlled mechanism.

The project's full classification fits in three named mechanisms,
collectively the \emph{master Mellin engine}~\cite{loutey_paper18,
loutey_paper32, loutey_paper55}:

\begin{center}
\begin{tabular}{@{}cllc@{}}
\toprule
Slot $k$ & Mechanism & Period-ring home & Source \\
\midrule
$k = 0$ & M1 (Hopf-base measure) & $\Q[\pi, \pi^{-1}]$ & Paper~55 \S 3 \\
$k = 2$ & M2 (Seeley--DeWitt) & $\bigoplus_k \pi^{2k}\cdot\Q$ & Paper~55 \S 4 \\
$k = 1$ & M3 (vertex parity / Hurwitz / Dirichlet $L$) & $\MT(\Z[i, 1/2])$ at level $\le 4$ & Paper~55 \S 5 \\
\bottomrule
\end{tabular}
\end{center}

The three rows are the three sub-cases of a single Mellin transform
\begin{equation}
\label{eq:master_mellin}
\mathcal{M}\!\left[\Tr\bigl(D^k \cdot e^{-tD^2}\bigr)\right](s)
\end{equation}
on the discrete Camporesi--Higuchi $\sthree$ spectral triple, with
operator-order index $k \in \{0, 1, 2\}$ selecting which mechanism
produces which kind of transcendental.

\paragraph{The classification is complete.}
Paper~32~\S~VIII~\cite{loutey_paper32} contains a
\emph{case-exhaustion theorem}:\ every appearance of $\pi$ in any
closed-form GeoVac spectral observable arises from one of M1, M2, or
M3.  No fourth mechanism has been found despite specific searches.
Paper~55~\cite{loutey_paper55} promotes this from a partition theorem
to a period-theoretic classification:\ every transcendental observed
in the corpus sits in cyclotomic mixed-Tate periods at level $\le 4$
over $\Z[i, 1/2]$ (Kontsevich--Zagier~\cite{kontsevich_zagier2001},
Brown~\cite{brown2012}, Deligne~\cite{deligne2010},
Glanois~\cite{glanois2015}, Fathizadeh--Marcolli
\cite{fathizadeh_marcolli2016} are the relevant published
classifications;\ each piece of GeoVac's data is placed inside one of
their named period rings).

\paragraph{What was a zoo became a classified program.}
Three years ago the project produced transcendentals through ad-hoc
calculations and was uncertain whether different appearances of $\pi$
or $\zeta(2)$ were ``the same'' or ``different.''  Today every
transcendental observation has a named motivic home, indexed by which
slot of \eqref{eq:master_mellin} produced it.  The classification
makes the operational question --- ``what kind of transcendental can
appear from \emph{this} calculation?'' --- algorithmic.

\paragraph{A note on periods.}  A period in the sense of
Kontsevich--Zagier 2001~\cite{kontsevich_zagier2001} is a complex
number whose real and imaginary parts can be expressed as integrals of
algebraic functions over domains defined by polynomial inequalities
with rational coefficients.  $\pi$, $\log 2$, $\zeta(n)$ for $n \ge
2$, Catalan's constant $G = \beta(2)$, the Dirichlet $\beta(s)$ values
at integer $s$ --- all are periods.  The cyclotomic mixed-Tate
sub-class at level $\le 4$ over $\Z[i, 1/2]$ is, conjecturally, a
small and well-understood corner of the full period algebra.  GeoVac
sits in this corner.

% =====================================================================
\section{The Forced/Free Seam}
\label{sec:forced_free}
% =====================================================================

The master Mellin engine of \S\ref{sec:periods} classifies what GeoVac
\emph{generates}.  An equally important question is what GeoVac
\emph{does not} generate --- in particular, the empirical constants of
physics (the fine-structure constant $\alpha$, the Yukawa couplings,
the generation count $N_{\mathrm{gen}}$, the cosmological constant)
that must enter as external inputs.  The framework's most operationally
useful structural discipline is the clean separation between what
geometry forces and what measurement supplies.

\paragraph{What geometry forces (the structural skeleton).}

The following are all theorem-grade forced by the framework's
construction, independent of any empirical input:

\begin{itemize}
  \item The outer manifold is $\sthree$.  Forced by Bertrand's
    theorem~\cite{bertrand1873} (only $1/r$ and harmonic-oscillator
    central potentials produce closed bound orbits, of which $1/r$ is
    the Coulomb case) plus the $SO(4)$ hidden symmetry of the Coulomb
    problem (Paper~7).
  \item The Hopf bundle structure $\sthree \to S^2$ is forced by the
    maximal-torus reduction of $SU(2) = \sthree$
    (Paper~25)~\cite{loutey_paper25}.
  \item The Camporesi--Higuchi Dirac sector is forced by the
    parallelizability of $\sthree$ (Paper~32)~\cite{loutey_paper32}.
  \item The Wilson $SU(2)$ gauge structure is forced by the
    non-abelian generalization of the abelian Hopf structure
    (Paper~30)~\cite{loutey_paper30}.
  \item The Standard Model gauge group $U(1) \times SU(2) \times
    SU(3)$ truncated at $n \le 3$ in the complex-Hopf tower
    $S^{2n-1} \to SU(n)$ is forced by Bertrand $\times$ Hopf-tower
    structural saturation (CLAUDE.md~\S~1.7 WH4 and the May~2026
    Door~4 Sprint).
  \item The dimension of the inner Dirac moduli space is
    $\dim \mathcal{M}(D_F) = 128$ per generation, $8$ in the diagonal
    slice.  Bertrand-shaped Forced-Count Theorem (Paper~32~\S~VIII,
    June~2026).
\end{itemize}

\paragraph{What measurement supplies (calibration data).}

The following are all empirical inputs to the framework, falling
outside any structural derivation it offers:

\begin{itemize}
  \item The fine-structure constant $\alpha^{-1} \approx 137.036$.
    Three independent spectral homes for the framework's combination
    $K = \pi(B + F - \Delta)$ have been established (B as Casimir
    trace, $F = \zeta(2)$ as Fock-Dirichlet, $\Delta = 1/40$ as
    Dirac mode degeneracy), but the combination rule itself remains
    a numerical observation, not derived.  Paper~2 is in the
    Observations folder~\cite{loutey_paper2};\ twelve mechanisms
    have been eliminated for a single-principle derivation of $K$.
  \item The Yukawa coupling values.  A 162-cell PSLQ sweep against
    low-coefficient pure-Tate periods (Sprint Yukawa-PSLQ,
    2026-06-03) returned zero hits.  The Yukawa Dirichlet ring lives
    in Paper~18~\S~IV.6's sixth tier (parameter-tied Dirichlet,
    categorically disjoint from M1/M2/M3).
  \item The generation count $N_{\mathrm{gen}} = 3$.  The
    Hopf-tower-to-representation extension was ruled out as a
    shortcut for forcing $N_{\mathrm{gen}}$
    (Sprint Read~2, 2026-06-03):\ the ``3 algebra factors'' and the
    ``3 generations'' are different $3$s in the standard SM
    representation.  Multi-year wall.
  \item The cosmological constant.  The framework's CC formula
    $\Lambda_{\mathrm{cc}} = 6\phi(2)/\phi(1) \cdot \Lambda^2$
    inherits the standard $\sim 10^{120}$ fine-tuning problem
    (Paper~51 G7).
\end{itemize}

\paragraph{The seam is itself a theorem.}

What makes the forced/free separation operationally useful, rather
than rhetorical, is that the boundary between the two categories is
itself controlled by named theorems:

\begin{itemize}
  \item The Forced-Count Theorem (Paper~32~\S~VIII):\ the gauge structure
    of GeoVac is forced.
  \item The Yukawa Non-Selection Theorem (Sprint H1
    POSITIVE-THIN closure, 2026-05-31):\ the framework's
    almost-commutative extension admits a Higgs structure but does
    not autonomously select the Yukawa matrix.
  \item The $\eta$-Trivialization Theorem and AC-Factorization
    Theorem (Paper~18~\S~IV.6):\ the combined heat trace on the AC
    extension factorises cleanly into an outer pure-Tate factor and
    an inner Dirichlet factor in the Yukawa variables, with no
    mixing.
\end{itemize}

These three theorems collectively pin the seam:\ values are free,
relations are forced (CLAUDE.md~\S~1.7 WH4 closure).  This is the
project's most operationally useful structural reading.

\paragraph{The C-arc closure (June 2026).}  The forced/free seam
discipline above was reframed as a theorem-grade structural
characterization between v3.97.0 and v3.103.0 (June 2026).  Paper~57
catalogues sixty named entries across nine domains (35 forced, 1
admitted-not-forced, 24 calibration); a five-candidate principle
probe (Sprint~C2) found that ``packing-reachability'' (P5) classifies
the catalogue at 98.3\% accuracy.  Paper~32~\S~VIII now consolidates
eight theorem-grade non-selection results (Forced-Count, H1 Yukawa,
$N_{\mathrm{gen}}$, KO-dimension, single-cutoff spectral action,
spatial-composition wall, no-single-mechanism-$K$, and the cutoff
function as external test-function data) across four sectors:
inner-factor structural skeleton, multi-focal composition,
$\alpha$-combination, and gravity.  The forced/free seam is no
longer an empirical observation;\ it is a theorem-bound boundary.

% =====================================================================
\section{What's in the Corpus}
\label{sec:corpus}
% =====================================================================

The GeoVac corpus is $50$+ papers organised into six audience-targeted
groups (\texttt{papers/group1\_operator\_algebras/}, etc.) plus an
archive and a synthesis folder.  Each group is described in CLAUDE.md
\S~6 with detailed paper-by-paper coverage.  This section gives
one-line summaries.

\begin{itemize}
  \item \textbf{Group~1 --- Operator algebras / NCG} ($14$ papers).
    Math.OA arc.  Headline:\ WH1 PROVEN (Paper~38) with five-lemma
    Gromov--Hausdorff convergence;\ Lorentzian propinquity
    first-in-literature (Paper~45);\ Mondino--S\"amann bridge
    (Paper~48);\ OSLPLS strong-form bridge (Paper~49).  Closes
    four published-open questions of Connes--van Suijlekom
    2021~\cite{connes_vs2021}.
  \item \textbf{Group~2 --- Quantum chemistry} ($9$ papers).
    Natural-geometry hierarchy levels~1--5.  Headline:\ structural
    accuracy ceilings characterized via the multi-focal-composition
    wall;\ Paper~8--9 and FCI-M are guardrail papers documenting
    proven negative results.
  \item \textbf{Group~3 --- Foundations} ($7$ papers including this
    field guide).  Headline:\ Paper~0 (packing), Paper~7
    (Fock-projection conformal equivalence), Paper~22 (angular
    sparsity), Paper~24 (Bargmann--Segal $\sfive$ lattice), Paper~31
    (universal/Coulomb partition), Paper~55 (periods classification).
  \item \textbf{Group~4 --- Quantum computing} ($4$ papers).
    Qubit-encoding resource estimates.  Headline:\ $O(Q^{2.5})$
    composed Pauli scaling, $51\times$--$1712\times$ advantage over
    Gaussian baselines, $28$ molecules in the library (Paper~14).
  \item \textbf{Group~5 --- QED / gauge / SM} ($7$ papers).
    Headline:\ hydrogen Lamb shift at $-0.534\%$ one-loop closure
    (Paper~36);\ SU(2) Wilson lattice gauge structure with
    maximal-torus reduction to U(1) Paper~25 (Paper~30);\ QED on
    $\sthree$ with five theorems and $S_{\min}$ irreducible at 200
    dps (Paper~28).
  \item \textbf{Group~6 --- Precision observations} ($4$ papers).
    Headline:\ five Roothaan autopsies decomposing observables into
    projection-chain contributions (Paper~34);\ time as projection
    (Paper~35);\ universal $S_B$ scaling on sparse lattices
    (Paper~27).
\end{itemize}

The \texttt{papers/archive/} folder contains historical and
superseded drafts.  The \texttt{papers/synthesis/} folder contains
this field guide plus two group-level syntheses (Group~1 and
Group~3) drafted May~2026.

% =====================================================================
\section{What's Open}
\label{sec:open}
% =====================================================================

The project is honest about scope.  The following are named open
questions, in rough order of how much external infrastructure each
would require:

\begin{itemize}
  \item \textbf{The $\alpha$-derivation}.  Paper~2's combination rule
    $K = \pi(B + F - \Delta)$ matches the empirical $\alpha^{-1}$ at
    $8.8 \times 10^{-8}$ relative precision, but is classified as a
    numerical observation, not a derivation.  Twelve mechanisms
    eliminated for a single-principle derivation.  CLAUDE.md~\S~1.7
    WH5 reads ``$\alpha$ is a projection constant, not a derivable
    number.''
  \item \textbf{The Yukawa values}.  Sprint Yukawa-PSLQ (2026-06-03)
    ruled out the Yukawas as low-coefficient pure-Tate periods.  They
    are Class~1 calibration data and live in Paper~18~\S~IV.6's sixth
    tier.  Whether they admit a deeper non-period derivation is open
    and falls outside the master Mellin engine's classification scope.
  \item \textbf{$N_{\mathrm{gen}} = 3$}.  Multi-year wall;\ no
    sprint-scale handle (Sprint Read~2, 2026-06-03).
  \item \textbf{The cosmological constant problem}.  Standard fine-tuning
    $\sim 10^{120}$ inherited from the Connes--Chamseddine spectral
    action expansion (Paper~51 G7).  No GeoVac-specific resolution
    proposed.
  \item \textbf{The chemistry W1e wall}.  Multi-determinant FCI correlation
    in second-row and heavier diatomic binding is below the
    framework's structural accuracy ceiling.  Six diagnostic sprints
    (May~2026) identified the wall as multi-focal-composition rather
    than a missing single mechanism.
  \item \textbf{The M2/M3 motivic Galois unification (Q5')}.  Both M2 and
    M3 engage the cyclotomic field $\Q(i) = \Q(\zeta_4)$, but via
    structurally independent mechanisms.  Whether a single Tannakian
    symmetry unifies them is an open mathematical research project
    (Sprint A7, 2026-06-03).  Multi-year.
  \item \textbf{Deligne--Milne pro-finite reconstruction (TC-2 converse)}.
    Paper~56 closes the cosmic-Galois inclusion $\Phi:\
    U^{*(n_{\max})} \hookrightarrow \mathrm{Aut}^\otimes(\omega)$ at
    finite cutoff via the Tannakian dual of an explicit Hopf algebra.
    The converse equality $\mathrm{Aut}^\otimes(\omega) =
    \lim_\leftarrow U^{*(n_{\max})}$ at the pro-finite limit
    (Deligne--Milne 1982 Theorem~2.11) is the named multi-year
    frontier;\ each TC-stone closed at finite cutoff (TC-1a..f, TC-2a..d)
    has compressed multi-year content into days, suggesting the
    pro-finite assembly may close at sprint-scale once the right
    transition-map machinery is in place.
  \item \textbf{Non-commutative Mondino--S\"amann (Q2')}.  Paper~49 closes
    the strong-form Krein--MS bridge on the enlarged-substrate
    chirality-flipping generators.  The non-commutative analog of
    Mondino--S\"amann pre-length-spaces (a non-commutative Lorentzian
    GH program built from spectral-triple data rather than causal
    diamonds) is the remaining multi-month NCG-research target.
  \item \textbf{Inhomogeneous extensions}.  The full
    Fathizadeh--Marcolli generic Robertson--Walker spacetime with
    non-trivial scaling factor $a(t)$ requires either a continuum
    substrate or a time-dependent Camporesi--Higuchi spectral triple,
    outside the discrete-substrate principle.  The closest GeoVac
    analog is the discrete-warped-substrate program of
    Paper~51~\S G4-3 (multi-month follow-on).
  \item \textbf{Full Standard Model spectral triple}.  G4a (Connes
    Standard Model on $\sthree$ alone) was closed POSITIVE-THIN in
    May~2026;\ G4b (genuine cross-manifold $\sthree \otimes
    \mathrm{Hardy}(\sfive)$) is blocked at the
    Coulomb-versus-harmonic-oscillator asymmetry of Paper~24~\S~V at
    layer~4 (modular Hamiltonian of the wedge KMS state).  Resolution
    requires NCG-framework extension.
\end{itemize}

What is \emph{not} open --- but is sometimes asked --- is whether
GeoVac is a new foundation for quantum mechanics.  It is not.  GeoVac
is a discretization framework whose discrete substrate is conformally
equivalent to known continuous quantum mechanics at every level where
the equivalence has been tested.  The math supports both readings;\
the interpretive question is left to the reader.

% =====================================================================
\section{Reader's Map}
\label{sec:readers_map}
% =====================================================================

This is a multi-paper corpus.  Here is where each kind of reader
should start.

\begin{itemize}
  \item \textbf{The mathematician (NCG, math.OA, spectral triples)}.
    Read Papers~38, 32, 45, 46, 47, 48, and 49.  Paper~38 is WH1
    PROVEN (SU(2) propinquity convergence on the Camporesi--Higuchi
    spectral triple);\ Paper~32 is the explicit spectral triple
    construction with Connes axiom audit and the \S~VIII C-arc
    closure of eight theorem-grade non-selection results;\
    Papers~45 and 46 are the first published Lorentzian propinquity
    convergence theorems (K$^+$-weak-form and strong-form,
    respectively) in the math.OA literature;\ Paper~47 closes the
    de-compactification limit at norm-resolvent level;\ Papers~48
    and 49 are the Mondino--S\"amann bridge and the OSLPLS
    strong-form bridge.
  \item \textbf{The mathematician (periods, motives, Tannakian)}.
    Read Papers~55 and 56.  Paper~55 places GeoVac's periods in the
    cyclotomic mixed-Tate ring at level $\le 4$ with explicit
    Glanois descent on the vertex-parity Mellin slot.  Paper~56
    identifies the cosmic-Galois group $U^* = \mathbb{G}_a^{3N}
    \rtimes SL_2$ at finite cutoff via the Tannakian dual of an
    explicit Hopf algebra $\mathcal{H}_{\mathrm{GV}}$, with 5{,}864
    bit-exact zero residuals across the Pro-System + Tannakian
    closure panel.  Inclusion direction closed at finite cutoff;
    converse equality with the Deligne--Milne reconstruction at
    pro-finite limit is the named multi-year frontier.  Open
    M2/M3 unification question (Q5') discussed in Paper~55~\S~7.6
    and Paper~56~\S sec:open\_g4.
  \item \textbf{The physicist (atomic, molecular, optical)}.  Read
    Papers~7, 22, and 14.  Paper~7 is the Fock-projection conformal
    equivalence;\ Paper~22 is the angular sparsity theorem (universal
    across central potentials);\ Paper~14 is the quantum-computing
    resource benchmarks.
  \item \textbf{The physicist (high-energy, gauge, SM)}.  Read
    Papers~30, 25, 28, and 36.  Paper~30 is SU(2) Wilson on
    $\sthree$;\ Paper~25 is the abelian-Hopf reduction;\ Paper~28 is
    QED on $\sthree$;\ Paper~36 is the bound-state QED arc closing
    the Lamb shift at $-0.534\%$ one-loop.
  \item \textbf{The physicist (gravity)}.  Read Paper~51.  Spectral
    action approach, two-term exact on $\sthree$, $S_{\mathrm{BH}} =
    A/4$ derivation pending the L6 sub-sprint (closed in May~2026).
  \item \textbf{The quantum computing practitioner}.  Read Paper~14
    and Paper~20.  Pauli scaling, OpenFermion/Qiskit/PennyLane
    export, market test against published Gaussian baselines.
  \item \textbf{The chemist}.  Read Papers~11, 13, 15, 17.  The
    natural-geometry hierarchy, levels~2--5.  Note the guardrail
    papers~8--9 and FCI-M before proposing any single-center or
    LCAO molecular encoding (these document proven negative
    results).
\end{itemize}

The CLAUDE.md document in the repository root is the project's
internal coordination file;\ it indexes every paper, every sprint,
every working hypothesis, and every documented failed approach.
External readers do not need to read it but may find the
``Failed approaches'' table (\S~3) instructive for understanding the
boundary between what works and what doesn't.

% =====================================================================
\section{Conclusion}
\label{sec:conclusion}
% =====================================================================

The one-sentence reading:\ \emph{GeoVac is a discrete spectral triple
whose periods classify the projective content of physics, sharply
separating what geometry forces from what measurement supplies.}

The story is a journey from a packing puzzle through Fock~1935's
$\sthree$ identification, through the natural-geometry hierarchy and
its chemistry investigation, to a structural reading as a discrete
spectral triple that instances Marcolli--van Suijlekom 2014 at the
data level and extends it via the forced-graph claim and master
Mellin engine, with
transcendentals classified as cyclotomic mixed-Tate periods at level
$\le 4$, and with the forced-versus-free seam pinned by three named
theorems.

The discipline is the separation:\ what geometry forces gets a name,
a theorem, and a paper;\ what measurement supplies is recorded as
calibration data and stays calibration data unless the framework can
produce a derivation, which so far it cannot.

The project began as an amateur geometric puzzle and has acquired the
status of a discrete spectral triple capable of closing
published-open questions in the math.OA literature.  That trajectory
was AI-augmented at every step:\ implementation, sprint memos,
literature collation, paper drafting, and the internal coordination
machinery were all developed collaboratively with Anthropic's Claude,
with PI direction and quality control.

The invitation:\ come read.  The corpus is dense, but the entry points
are now indexed.  Section~\ref{sec:readers_map} above is the
reader's map.

% =====================================================================
\paragraph*{Acknowledgments.}
This work was developed in an AI-augmented agentic workflow with
Anthropic's Claude.  Implementation, internal sprint memos, paper
drafting, and bibliographic collation were performed collaboratively;\
all theorems, design choices, and editorial decisions are the
author's responsibility.  The project has no institutional affiliation
and is disseminated through DOI-stamped Zenodo releases.

% =====================================================================
\begin{thebibliography}{99}

\bibitem{fock1935}
V.~Fock, ``Zur Theorie des Wasserstoffatoms,''
\textit{Z. Physik}~\textbf{98}, 145--154 (1935).

\bibitem{bertrand1873}
J.~Bertrand, ``Th\'eor\`eme relatif au mouvement d'un point attir\'e
vers un centre fixe,''
\textit{C.~R.\ Acad.\ Sci.\ Paris}~\textbf{77}, 849--853 (1873).

\bibitem{connes1994}
A.~Connes, \textit{Noncommutative Geometry}, Academic Press (1994).

\bibitem{connes_vs2021}
A.~Connes and W.~D.~van Suijlekom, ``Spectral truncations in
noncommutative geometry and operator systems,''
\textit{Commun.\ Math.\ Phys.}~\textbf{383}, 2021--2059 (2021).

\bibitem{marcolli_vs2014}
M.~Marcolli and W.~D.~van Suijlekom, ``Gauge networks in
noncommutative geometry,''
\textit{J.\ Geom.\ Phys.}~\textbf{75}, 71--91 (2014);\ preprint
arXiv:1301.3480 (2013).

\bibitem{latremoliere2017}
F.~Latr\'emoli\`ere, ``The dual Gromov-Hausdorff propinquity,''
\textit{J.\ Math.\ Pures Appl.}~\textbf{103}, 303--351 (2015).

\bibitem{kontsevich_zagier2001}
M.~Kontsevich and D.~Zagier, ``Periods,'' in \textit{Mathematics
Unlimited --- 2001 and Beyond}, Springer (2001), pp.~771--808.

\bibitem{brown2012}
F.~Brown, ``Mixed Tate motives over $\Z$,''
\textit{Ann.\ of Math.}~(2)~\textbf{175}, 949--976 (2012).

\bibitem{deligne2010}
P.~Deligne, ``Le groupe fondamental unipotent motivique de
$\mathbb{G}_m - \mu_N$, pour $N = 2, 3, 4, 6$ ou $8$,''
\textit{Publ.\ Math.\ IHES}~\textbf{112}, 101--141 (2010).

\bibitem{glanois2015}
C.~Glanois, ``Motivic unipotent fundamental groupoid of
$\mathbb{G}_m \setminus \mu_N$ for $N = 2, 3, 4, 6, 8$ and Galois
descents,'' arXiv:1411.4947 (2015);\ \textit{J.\ Number Theory}~\textbf{182},
36--90 (2018).

\bibitem{fathizadeh_marcolli2016}
F.~Fathizadeh and M.~Marcolli, ``Periods and motives in the spectral
action of Robertson--Walker spacetimes,''
\textit{Commun.\ Math.\ Phys.}~\textbf{356}, 641--671 (2017);\ preprint
arXiv:1611.01815 (2016).

\bibitem{loutey_paper0}
J.~Loutey, ``Paper 0: Geometric Packing and the Universal Constant
$\kappa = -1/16$,'' GeoVac corpus (2026).

\bibitem{loutey_paper2}
J.~Loutey, ``Paper 2: $\alpha = K^{-1}$ as a spectral coincidence on
the Fock-projected $\sthree$,'' GeoVac corpus (2026).

\bibitem{loutey_paper7}
J.~Loutey, ``Paper 7: The Dimensionless Vacuum on $\sthree$,''
GeoVac corpus (2026).

\bibitem{loutey_paper14}
J.~Loutey, ``Paper 14: Qubit encoding, Pauli scaling, and composed
sparsity,'' GeoVac corpus (2026).

\bibitem{loutey_paper18}
J.~Loutey, ``Paper 18: Exchange constants and the master Mellin
engine,'' GeoVac corpus (2026).

\bibitem{loutey_paper25}
J.~Loutey, ``Paper 25: GeoVac Hopf graphs as discrete lattice-gauge
structure,'' GeoVac corpus (2026).

\bibitem{loutey_paper30}
J.~Loutey, ``Paper 30: SU(2) Wilson lattice gauge on $\sthree =
SU(2)$,'' GeoVac corpus (2026).

\bibitem{loutey_paper32}
J.~Loutey, ``Paper 32: The GeoVac Spectral Triple,''
GeoVac corpus (2026).

\bibitem{loutey_paper38}
J.~Loutey, ``Paper 38: SU(2) propinquity convergence of truncated
Camporesi--Higuchi spectral triples,'' GeoVac corpus (2026);\
released to Zenodo 2026-05-07.

\bibitem{loutey_paper40}
J.~Loutey, ``Paper 40: Universal $4/\pi$ rate constant in
Latr\'emoli\`ere propinquity across compact connected Lie groups,''
GeoVac corpus (2026).

\bibitem{loutey_paper42}
J.~Loutey, ``Paper 42: Tomita--Takesaki modular structure and the
four-witness Wick-rotation theorem,'' GeoVac corpus (2026).

\bibitem{loutey_paper43}
J.~Loutey, ``Paper 43: Lorentzian extension at finite cutoff via
Krein-space at signature (3, 1),'' GeoVac corpus (2026).

\bibitem{loutey_paper44}
J.~Loutey, ``Paper 44: Operator-system extension of the BBB Krein
spectral triple,'' GeoVac corpus (2026).

\bibitem{loutey_paper45}
J.~Loutey, ``Paper 45: K$^+$-restricted weak-form Lorentzian
propinquity convergence,'' GeoVac corpus (2026).

\bibitem{loutey_paper46}
J.~Loutey, ``Paper 46: Strong-form Lorentzian propinquity on the
enlarged substrate,'' GeoVac corpus (2026).

\bibitem{loutey_paper47}
J.~Loutey, ``Paper 47: Norm-resolvent convergence to the non-compact
Lorentzian Krein spectral triple,'' GeoVac corpus (2026).

\bibitem{loutey_paper48}
J.~Loutey, ``Paper 48: Krein-pointed proper QMS and the bridge to
Mondino--S\"amann Lorentzian pre-length spaces,''
GeoVac corpus (2026).

\bibitem{loutey_paper49}
J.~Loutey, ``Paper 49: OSLPLS strong-form bridge and the
twin-paradox structural identification,'' GeoVac corpus (2026).

\bibitem{loutey_paper55}
J.~Loutey, ``Paper 55: Periods of GeoVac,'' GeoVac corpus (2026).

\end{thebibliography}

\end{document}
